\section{Experimental Setup}
\label{sec:setup}

\begin{figure*}[t]
\centering
\begin{tikzpicture}[
    font=\scriptsize,
    box/.style={draw,rounded corners=1.5pt,minimum height=8mm,minimum width=19mm,align=center,inner sep=2.5pt,fill=black!4},
    sec/.style={box,fill=teal!14,draw=teal!55!black},
    ar/.style={-{Stealth[length=4pt]},draw=black!60}
  ]
  \node[sec] (splits) {frozen splits\\{\tiny opt / sel / final, hashed}};
  \node[sec,right=4.5mm of splits] (iso) {isolation\\{\tiny agent view, uid, mount;}\\{\tiny negative test}};
  \node[box,right=4.5mm of iso] (search) {search\\{\tiny one optimizer,}\\{\tiny 3 seeds per arm}};
  \node[sec,right=4.5mm of search] (rule) {promotion\\{\tiny rule P or native}\\{\tiny rule per arm}};
  \node[sec,right=4.5mm of rule] (bstar) {matched\\{\tiny checkpoint $N^\star$}\\{\tiny from counters only}};
  \node[sec,right=4.5mm of bstar] (final) {final set\\{\tiny not used by search,}\\{\tiny scored once}};
  \node[sec,right=4.5mm of final] (stats) {statistics\\{\tiny paired, family-}\\{\tiny clustered, Holm-adjusted}};
  \foreach \a/\b in {splits/iso,iso/search,search/rule,rule/bstar,bstar/final,final/stats}
    \draw[ar] (\a) -- (\b);
  \draw[ar] (rule.north) to[bend right=35] node[above,font=\tiny] {reject / promote, next candidate} (search.north);
  \fill[teal!25,draw=teal!55!black] (0,-1.1) rectangle ++(0.26,0.13);
  \node[anchor=west,font=\tiny] at (0.34,-1.04) {fixed by the protocol before any outcome};
  \fill[black!4,draw=black!60] (6.6,-1.1) rectangle ++(0.26,0.13);
  \node[anchor=west,font=\tiny] at (6.94,-1.04) {the method under test};
\end{tikzpicture}
\caption{The measurement pipeline. Everything shaded teal is fixed before any search outcome is seen and is identical across arms, except that competitors keep their native acceptance rule. Search sees only the optimization and selection splits. The final set is scored once per frozen artifact.}
\label{fig:pipeline}
\end{figure*}

\subsection{Suites, models and splits}

Table~\ref{tab:suites} lists the three suites. Suite A edits spreadsheets with a code-execution agent: the target is Gemma-4-31B-it, served by one IBM RITS (Research Inference Service) deployment, on the 400-task verified subset of the 912 SpreadsheetBench questions \citep{spreadsheetbench}\footnote{File \texttt{spreadsheetbench\_verified\_400.tar.gz} (sha256 \texttt{10ef893d}\ldots) in \texttt{github.com/RUCKBReasoning/SpreadsheetBench}, \texttt{data/} directory.}. Suites B and C are customer-service dialogue in $\tau^2$-bench retail and airline \citep{tau2-bench,tau_bench}, pinned at commit 8ebb7499, with gpt-oss-120b as target and user simulator. Suites B and C are descriptive: with 40 and 20 final tasks and a simulated user they cannot resolve small differences, so they show that the protocol transfers, and inference there is at the artifact level only. GPT-5.5 grades natural-language assertions for every arm, because the benchmark's default judge is not served by our gateway and the target should not grade itself. Claude Sonnet 5.5 is the optimizer or reflection model for every matched arm. One metering gateway tags every model call with suite, arm, search seed and role.

\begin{table}[t]
\caption{Suites. Opt is the tasks whose rollouts the optimizer reads; sel is the tasks that gate acceptance; final is scored once per artifact. Final size is tasks $\times$ trials per task. *Sealed against scoring and selection; the task files were filesystem-visible to the optimizer, and no access was observed. Stock \capevolve{} reads and gates on one validation set. Suite A's final set is a stratified sample of 239 ids from the 912-question release, each scored at most twice before (score only).}
\label{tab:suites}
\centering\small
\begin{tabularx}{\columnwidth}{@{}lXlll@{}}
\toprule
 & Target & Opt / sel & Final & Exposure \\
\midrule
A & Gemma-4-31B-it & 80 / 40 & $70\times 1$ & score-only exposed \\
B & gpt-oss-120b & 49 / 25 & $40\times 3$ & sealed* \\
C & gpt-oss-120b & 20 / 10 & $20\times 3$ & development-exposed \\
\bottomrule
\end{tabularx}
\end{table}

On Suite A the reward is the LibreOffice-recalculated hard score; the as-saved score is logged for every rollout. The SkillOpt split (80 train, 40 validation) supplies optimization and selection. Competitors read train traces and select on validation. The stock \capevolve{} algorithms read validation traces and gate on validation, so they see 40 tasks where competitors see 120. This differs in exposure, in a direction we cannot sign, and we keep it as the native design. Proposers are isolated by user id or mount namespace with a negative test before launch, except where Appendix~\ref{app:incidents} records a gap: on Suite B, agent-optimize can read gold actions of its own optimization and selection tasks. Search-time selection scores are not comparable across methods, since GEPA selects its best on the same selection tasks it reports, so only final sets are compared. Every arm receives the same scorer feedback, which carries gold-derived facts such as fill counts and error direction but no gold values.

The Suite A final set is a stratified sample from the 239 ids of the 912-question release that remain after removing development-exposed families. Each id was scored at most in two earlier runs, score only, never used for design. We fix the sampling rule and seed, and hash the sample, before any final-set rollout is scored.

\subsection{Arms}

Table~\ref{tab:arms} lists the arms; an arm is one method with its configuration, and the seed artifact is the unoptimized starting point of every arm. Eight of the nine competitors run from official repositories at pinned commits (Appendix~\ref{app:config}); WikiSkill has no official code, so we run a community port. VERSE, CoEvoSkills, and HarnessLens on Suite A failed to reproduce (Appendix~\ref{app:incidents}). Five more methods (AHE, MCE, RHO, Memento-Skills, GraphSkillEvo) are one-seed exploratory runs, fixed before any outcome, outside the headline tables. Eight competing methods run on Suite A, GEPA on Suites B and C, and HarnessLens on Suite B.

\begin{table}[t]
\caption{Arms. Main arms run three search seeds to the matched checkpoint; ``cheap'' arms run three search seeds with a native budget that is usually below the cap and report actual rollouts. Pins are in Appendix~\ref{app:config}.}
\label{tab:arms}
\centering\small
\begin{tabularx}{\columnwidth}{@{}lXc@{}}
\toprule
Arm & Notes & Suites \\
\midrule
\capevolve{} hill-climb & rule P, last promoted candidate & A B C \\
\capevolve{} agent-optimize & agent-driven engine, rule P & A B C \\
GEPA \citep{gepa} & v0.1.4, Pareto reflective search & A B C \\
SkillOpt \citep{skillopt} & train feedback, val selection & A \\
Meta-Harness \citep{meta-harness} & container-isolated proposer & A \\
EvoSkill \citep{evoskill} & official loop & A \\
WikiSkill \citep{wikiskill} & community port, no official code & A \\
Trace2Skill \citep{trace2skill} & cheap, about 200 rollouts & A \\
ACE \citep{ace} & cheap, field renamed (App.~\ref{app:incidents}) & A \\
MIPROv2 \citep{miprov2-dspy} & cheap & A \\
HarnessLens \citep{xu2026harnesslens} & native tau2 fit & B \\
Unoptimized seed & one shared measurement per suite & A B C \\
One-shot repair & one optimizer step, accepted unconditionally & B C \\
\bottomrule
\end{tabularx}
\end{table}

\subsection{Matched checkpoint and acceptance rules}

Each search may count a budget of $N=600$ target rollouts, every evaluation included, and up to 4M non-cache-read optimizer tokens. One Gemma deployment gives about 12 rollouts per hour per search, so slow arms miss 600. We therefore compare at a matched checkpoint $N^\star$: the largest rollout count that every main seed of every main arm reaches by the search cutoff (2026-10-13 06:00 UTC), rounded down to a multiple of 50 and floored at 200. $N^\star$ comes from rollout counters of the primary-family arms only; the other arms are compared at the smaller of $N^\star$ and their own count. A candidate whose evaluation straddles $N^\star$ is excluded. Each arm is scored on what its own selection rule would have deployed at $N^\star$, using only information produced inside its first $N^\star$ rollouts. A seed below 200 rollouts is reported as incomplete and its arm leaves the table. The placeholder values are $N^\star=$ \mockBstarA{}, \mockBstarB{} and \mockBstarC{} for Suites A, B and C. \needs{Final $N^\star$ and the list of incomplete seeds, after the search cutoff.}

The matching is on target rollouts. It is not on optimizer compute, which differs by arm, so we report optimizer tokens and dollars per arm and make no claim of compute parity. It is also not on acceptance rules: \capevolve{} judges candidates with rule P and competitors keep their native rules. We therefore compare each method together with its own acceptance rule, and treat rule P as a deliberately conservative configuration. Exploratory stock-gate hill-climb arms on Suites A and C (three search seeds) show what the rule costs and buys; we compare them with rule-P hill-climb at the same rollout count, descriptively, without p-values. For every arm and seed we report candidates screened, candidates promoted under the native rule, and whether the returned artifact equals the seed artifact. Truncation at $N^\star$ can penalise arms whose schedule depends on the full budget (SkillOpt's cosine schedule, EvoSkill's iteration cap), so a secondary analysis scores each arm's native end-of-run artifact at the cutoff, labelled as unmatched.

\subsection{Statistics}
\label{sec:setup:stats}

The unit of inference is the task, clustered by task family. We separate three quantities: the score on the fixed final set (execution noise only), the expected score on the population the final set was drawn from (family-clustered interval), and the spread of final artifacts across search seeds. The headline is the paired population-level contrast, computed with cluster-robust standard errors and a $t$ interval, with Holm's step-down adjustment for multiple comparisons inside each suite. The primary family pairs each \capevolve{} algorithm (hill-climb, agent-optimize) with the seed artifact and the main competitors: 12 contrasts on Suite A, 8 on B and 6 on C. Appendix~\ref{app:sim} lists them, and the secondary families with their own Holm corrections.

Each primary row carries two intervals. The artifact-level interval, used for tests and the Holm correction, supports statements about the three frozen artifacts as found in these runs. Any statement about the method needs the pipeline-level interval ($t$ with two degrees of freedom, adding seed variance); if it includes zero, we restrict the statement to the frozen artifacts. Per-seed values accompany every mean. Every infrastructure error is retried up to three times with new trial seeds, and a cell still failing counts as a failure (score 0) for every arm and the seed artifact. Hangs past three hours get one retry and then count as missing. A win needs a Holm-adjusted $p<0.05$ and an interval that excludes zero; a sensitivity analysis treats failed cells as missing, with pessimistic and optimistic bounds, and we report both if they disagree. Two artifacts are equivalent only if two one-sided tests at $\alpha=0.05$ both reject for a margin of $0.10$ (Suites A and B); otherwise the result is inconclusive, and we do not write ``no difference''.

Variance comes from the seed artifact's observed evaluations, not assumptions (Appendix~\ref{app:sim}). The minimum detectable effect (MDE) for contrasts between the three frozen artifacts, at 80\% power and $\alpha=0.05$ before Holm, is \mdeACondOne{} points on Suite A (70 tasks, one trial), \mdeBCondThree{} on B (40 tasks, three trials) and \mdeCCondThree{} on C (20 tasks, three trials); claims about the optimizer itself need \mdeAPipeOne{}, \mdeBPipeThree{} and \mdeCPipeThree{} points. Suite A's seed-artifact mean moved by 8.8 points across six evaluation batches, so we score all arms interleaved in one window, with the seed artifact scored twice. 

\subsection{Cost accounting}

Every call is priced from the gateway log. Claude models use list prices as a proxy; RITS and Azure-served models have no visible price, so their tokens are reported and never counted as zero, and totals that include them are lower bounds (Appendix~\ref{app:cost}).
